% Bericht zu Mensch ärgere dich nicht: EvE- vs. PvE-Statistik.
% Vorher: python collect_stats.py   (erzeugt generated/*.csv und generated/pies.tex)
% Dann:   pdflatex report.tex  (zweimal nicht nötig)
\documentclass[a4paper,11pt]{article}
\usepackage[utf8]{inputenc}
\usepackage[T1]{fontenc}
\usepackage[ngerman]{babel}
\usepackage[margin=2.2cm]{geometry}
\usepackage{amsmath}
\usepackage{pgfplots}
\usepackage{pgf-pie}
\usepackage{float}
\pgfplotsset{compat=1.16, table/col sep=comma}

\def\modelist{{Classic 4 Player Field},{Classic 6 Player Field},{2 Players on classic 6 Player Field},{2 Players on classic 4 Player Field},{3 Players on classic 6 Player Field},{2 Players on 24 Tiles},{3 Players on 30 Tiles},{3 Players on classic 4 Player Field}}
\def\modelast{7}

\definecolor{eveColor}{RGB}{200,40,40}
\definecolor{pveColor}{RGB}{35,90,200}

% #1 = Metrik, #2 = Achsenbeschriftung, #3 = ymax (Platz für die Beschriftung über den Balken)
\newcommand{\lengthplot}[3]{%
\begin{tikzpicture}
\begin{axis}[
  bar width=10pt,
  width=\textwidth, height=10cm,
  ymin=0, ymax=#3,
  ylabel={#2},
  xtick={0,...,\modelast},
  xticklabels/.expand once=\modelist,
  xticklabel style={rotate=90, anchor=east, font=\small},
  enlarge x limits=0.07,
  ymajorgrids, major grid style={dotted,gray},
  legend style={at={(0.5,1.04)}, anchor=south, legend columns=-1, font=\small},
  error bars/error bar style={black!70},
  error bars/error mark options={mark size=3pt, black!70},
]
\addplot+[ybar, fill=eveColor, draw=eveColor!70!black, bar shift=0pt, mark=none,
  error bars/.cd, y dir=both, y explicit]
  table[x expr=\coordindex-0.2, y=EvE_#1_mean, y error=EvE_#1_std] {generated/lengths.csv};
\addplot+[ybar, fill=pveColor, draw=pveColor!70!black, bar shift=0pt, mark=none,
  error bars/.cd, y dir=both, y explicit]
  table[x expr=\coordindex+0.2, y=PvE_#1_mean, y error=PvE_#1_std] {generated/lengths.csv};
% unsichtbare Balken auf Höhe Mittelwert + StabW tragen die Beschriftung (Wert und n)
\addplot[ybar, draw=none, fill=none, bar shift=0pt, mark=none, forget plot,
  nodes near coords, point meta=explicit symbolic,
  every node near coord/.style={rotate=90, anchor=west, font=\scriptsize, text=eveColor!80!black}]
  table[x expr=\coordindex-0.2, y=EvE_#1_top, meta=EvE_#1_label] {generated/lengths.csv};
\addplot[ybar, draw=none, fill=none, bar shift=0pt, mark=none, forget plot,
  nodes near coords, point meta=explicit symbolic,
  every node near coord/.style={rotate=90, anchor=west, font=\scriptsize, text=pveColor!80!black}]
  table[x expr=\coordindex+0.2, y=PvE_#1_top, meta=PvE_#1_label] {generated/lengths.csv};
\legend{EvE (Computer gegen Computer), PvE (Mensch gegen Computer)}
\end{axis}
\end{tikzpicture}}
\title{Mensch ärgere dich nicht\\Statistischer Vergleich von EvE- und PvE-Spielen}
\author{}
\date{\today}

\begin{document}
\maketitle

\section{Datengrundlage}
Ausgewertet werden alle gespeicherten Spielprotokolle je Spielmodus aus
\texttt{game\_mode.json}. \textbf{EvE}-Spiele (rot) sind simulierte Partien, in denen
alle Spieler zufällig unter den erlaubten Zügen wählen; \textbf{PvE}-Spiele (blau) sind
Partien, in denen ein Mensch gegen zufällig ziehende Computerspieler antritt. Die
Stichprobengröße $n$ steht an jedem Balken. Modi ohne PvE-Spiele zeigen $n=0$ und
einen leeren Balken.

Einige Modi nutzen ein Brett, das für mehr Spieler ausgelegt ist, als tatsächlich
mitspielen (z.\,B.\ drei Spieler auf dem Vierer-Brett, Sitze 1, 2 und 4). Dadurch sind
die Abstände zwischen den Startfeldern ungleich.

\section{Spiellänge}
Alle drei Diagramme zeigen den Mittelwert je Spielmodus. Der schwarze Fehlerbalken
ist die Standardabweichung der Spiellänge. Je Modus stehen zwei Balken nebeneinander:
oben EvE (rot), darunter PvE (blau).

\subsection{Würfe}
Gesamtzahl der Würfelwürfe bis zum Spielende, einschließlich der bis zu drei Versuche,
mit denen ein Spieler versucht, eine Figur aus dem Haus zu bekommen.

\begin{figure}[H]\centering
\lengthplot{rolls}{Würfe pro Spiel}{3000}
\caption{Durchschnittliche Anzahl der Würfelwürfe pro Spiel.}
\end{figure}

\subsection{Runden}
Eine Runde ist ein vollständiger Umlauf, in dem jeder Spieler einmal an der Reihe war.
Spiele mit weniger Spielern oder kürzeren Brettern enden in weniger Würfen, nicht
unbedingt in weniger Runden.

\begin{figure}[H]\centering
\lengthplot{rounds}{Runden pro Spiel}{550}
\caption{Durchschnittliche Anzahl der Runden pro Spiel.}
\end{figure}

\subsection{Einzelne Züge}
Gezählt werden alle einzelnen Zugeinträge im Protokoll. Würfelt ein Spieler eine
Sechs, zieht er erneut, was einen zusätzlichen Zug erzeugt. Die Zahl liegt daher über
der Rundenzahl multipliziert mit der Spielerzahl.

\begin{figure}[H]\centering
\lengthplot{moves}{Züge pro Spiel}{2800}
\caption{Durchschnittliche Anzahl einzelner Züge pro Spiel.}
\end{figure}

\section{Zugcharakteristik}
\subsection{Würfe pro Zug}
Mittlere Zahl der Würfelwürfe je Zug (Zugeintrag). Der Wert liegt knapp über~1, weil
nur ein Spieler ohne Figur auf dem Brett bis zu dreimal würfeln darf, um eine Sechs zu
erhalten.

\begin{figure}[H]\centering
\lengthplot{rolls_per_move}{Würfe pro Zug}{1.8}
\caption{Durchschnittliche Zahl der Würfe pro Zug.}
\end{figure}

\subsection{Anteil der Züge mit Entscheidung}
Anteil der Züge, in denen der Spieler tatsächlich zwischen mehreren sinnvollen Figuren
wählen musste. Ein Wert von 30\,\% heißt: Nur etwa jeder dritte Zug erfordert eine
Entscheidung; in allen anderen Zügen ist der Zug erzwungen oder es gibt keine Wahl
(z.\,B.\ gilt die Hausregel, dass Einziehen oder Vorrücken im Haus Vorrang hat).

\begin{figure}[H]\centering
\lengthplot{decision_share}{Züge mit Entscheidung in \%}{50}
\caption{Anteil der Züge mit echter Entscheidungsmöglichkeit.}
\end{figure}

\section{Vorsprung des Siegers}
Beim Spielende wird gezählt, wie viele Figuren die übrigen Spieler bereits im Haus
hatten, gemittelt über alle Gegner. Ein kleiner Wert bedeutet einen deutlichen
Vorsprung des Siegers, ein hoher Wert (maximal die Figurenzahl) ein knappes Rennen.

\begin{figure}[H]\centering
\lengthplot{opponent_house}{Figuren der Gegner im Haus}{4}
\caption{Durchschnittliche Zahl der Figuren im Haus je Gegner bei Spielende.}
\end{figure}
\section{Siegverteilung}
Anteil der Siege je Spieler (Farben: Spieler~1 rot, 2 blau, 3 grün, 4 orange,
5 violett, 6 braun). Links steht jeweils EvE, rechts PvE. Bei fairem Spiel gewinnt
jeder Spieler mit dem Anteil $1/n_\text{Spieler}$. Beim PvE-Modus ist Spieler~1 der
Mensch.

\subsection*{Classic 4 Player Field}
{\small 4 Spieler; erwarteter Anteil je Spieler: $25.0\,\%$}\par
\begin{center}
\begin{tabular}{c@{\hspace{2cm}}c}
\textbf{EvE} & \textbf{PvE} \\
\begin{tabular}[t]{@{}c@{}}\begin{tikzpicture}\pie[radius=1.6, color={red,blue,green,orange}, text=legend, rotate=90, after number=\%]{26.5/1, 25.4/2, 23.8/3, 24.3/4}\end{tikzpicture}\\[0.2em] {\small $n=1000$}\end{tabular} &
\begin{tabular}[t]{@{}c@{}}\begin{tikzpicture}\pie[radius=1.6, color={red,blue,green,orange}, text=legend, rotate=90, after number=\%]{44.4/1, 18.5/2, 24.1/3, 13.0/4}\end{tikzpicture}\\[0.2em] {\small $n=54$}\end{tabular} \\
\end{tabular}
\end{center}

\subsection*{Classic 6 Player Field}
{\small 6 Spieler; erwarteter Anteil je Spieler: $16.7\,\%$}\par
\begin{center}
\begin{tabular}{c@{\hspace{2cm}}c}
\textbf{EvE} & \textbf{PvE} \\
\begin{tabular}[t]{@{}c@{}}\begin{tikzpicture}\pie[radius=1.6, color={red,blue,green,orange,purple,brown}, text=legend, rotate=90, after number=\%]{15.8/1, 17.9/2, 17.1/3, 16.9/4, 16.7/5, 15.6/6}\end{tikzpicture}\\[0.2em] {\small $n=1000$}\end{tabular} &
\begin{tabular}[t]{@{}c@{}}\begin{tikzpicture}\pie[radius=1.6, color={red,blue,green,orange,purple,brown}, text=legend, rotate=90, after number=\%]{46.7/1, 6.7/2, 13.3/3, 0.0/4, 20.0/5, 13.3/6}\end{tikzpicture}\\[0.2em] {\small $n=15$}\end{tabular} \\
\end{tabular}
\end{center}

\subsection*{2 Players on classic 6 Player Field}
{\small 2 Spieler; erwarteter Anteil je Spieler: $50.0\,\%$}\par
\begin{center}
\begin{tabular}{c@{\hspace{2cm}}c}
\textbf{EvE} & \textbf{PvE} \\
\begin{tabular}[t]{@{}c@{}}\begin{tikzpicture}\pie[radius=1.6, color={red,blue}, text=legend, rotate=90, after number=\%]{51.4/1, 48.6/2}\end{tikzpicture}\\[0.2em] {\small $n=1000$}\end{tabular} &
\begin{tabular}[t]{@{}c@{}}\emph{keine Daten}\end{tabular} \\
\end{tabular}
\end{center}

\subsection*{2 Players on classic 4 Player Field}
{\small 2 Spieler; erwarteter Anteil je Spieler: $50.0\,\%$}\par
\begin{center}
\begin{tabular}{c@{\hspace{2cm}}c}
\textbf{EvE} & \textbf{PvE} \\
\begin{tabular}[t]{@{}c@{}}\begin{tikzpicture}\pie[radius=1.6, color={red,blue}, text=legend, rotate=90, after number=\%]{51.4/1, 48.6/2}\end{tikzpicture}\\[0.2em] {\small $n=1000$}\end{tabular} &
\begin{tabular}[t]{@{}c@{}}\begin{tikzpicture}\pie[radius=1.6, color={red,blue}, text=legend, rotate=90, after number=\%]{50.0/1, 50.0/2}\end{tikzpicture}\\[0.2em] {\small $n=14$}\end{tabular} \\
\end{tabular}
\end{center}

\subsection*{3 Players on classic 6 Player Field}
{\small 3 Spieler; erwarteter Anteil je Spieler: $33.3\,\%$}\par
\begin{center}
\begin{tabular}{c@{\hspace{2cm}}c}
\textbf{EvE} & \textbf{PvE} \\
\begin{tabular}[t]{@{}c@{}}\begin{tikzpicture}\pie[radius=1.6, color={red,blue,green}, text=legend, rotate=90, after number=\%]{35.0/1, 32.7/2, 32.3/3}\end{tikzpicture}\\[0.2em] {\small $n=1000$}\end{tabular} &
\begin{tabular}[t]{@{}c@{}}\emph{keine Daten}\end{tabular} \\
\end{tabular}
\end{center}

\subsection*{2 Players on 24 Tiles}
{\small 2 Spieler; erwarteter Anteil je Spieler: $50.0\,\%$}\par
\begin{center}
\begin{tabular}{c@{\hspace{2cm}}c}
\textbf{EvE} & \textbf{PvE} \\
\begin{tabular}[t]{@{}c@{}}\begin{tikzpicture}\pie[radius=1.6, color={red,blue}, text=legend, rotate=90, after number=\%]{49.1/1, 50.9/2}\end{tikzpicture}\\[0.2em] {\small $n=1000$}\end{tabular} &
\begin{tabular}[t]{@{}c@{}}\begin{tikzpicture}\pie[radius=1.6, color={red,blue}, text=legend, rotate=90, after number=\%]{57.1/1, 42.9/2}\end{tikzpicture}\\[0.2em] {\small $n=42$}\end{tabular} \\
\end{tabular}
\end{center}

\subsection*{3 Players on 30 Tiles}
{\small 3 Spieler; erwarteter Anteil je Spieler: $33.3\,\%$}\par
\begin{center}
\begin{tabular}{c@{\hspace{2cm}}c}
\textbf{EvE} & \textbf{PvE} \\
\begin{tabular}[t]{@{}c@{}}\begin{tikzpicture}\pie[radius=1.6, color={red,blue,green}, text=legend, rotate=90, after number=\%]{34.2/1, 32.2/2, 33.6/3}\end{tikzpicture}\\[0.2em] {\small $n=1000$}\end{tabular} &
\begin{tabular}[t]{@{}c@{}}\emph{keine Daten}\end{tabular} \\
\end{tabular}
\end{center}

\subsection*{3 Players on classic 4 Player Field}
{\small 3 Spieler, Brett für 4, Sitze 1, 2, 4; erwarteter Anteil je Spieler: $33.3\,\%$}\par
\begin{center}
\begin{tabular}{c@{\hspace{2cm}}c}
\textbf{EvE} & \textbf{PvE} \\
\begin{tabular}[t]{@{}c@{}}\begin{tikzpicture}\pie[radius=1.6, color={red,blue,green}, text=legend, rotate=90, after number=\%]{36.0/1, 31.5/2, 32.5/3}\end{tikzpicture}\\[0.2em] {\small $n=1000$}\end{tabular} &
\begin{tabular}[t]{@{}c@{}}\begin{tikzpicture}\pie[radius=1.6, color={red,blue,green}, text=legend, rotate=90, after number=\%]{100.0/1, 0.0/2, 0.0/3}\end{tikzpicture}\\[0.2em] {\small $n=2$}\end{tabular} \\
\end{tabular}
\end{center}


\end{document}
